\subsection{Pollard's Rho (Factorización Rápida de Enteros)}

\graybold{Preguntas guía:}

\begin{itemize}
\item ¿Necesito factorizar un número entero muy grande (hasta aproximadamente \(10^{18}\))?
\item ¿La división por prueba o la criba son demasiado lentas debido al tamaño del número?
\item ¿Debo obtener rápidamente los factores primos para resolver un problema de teoría de números?
\end{itemize}

\graybold{Utilidad:}

Factorizar enteros grandes de forma probabilística utilizando la detección de ciclos sobre una función pseudoaleatoria. Combinado con Miller-Rabin para verificar primalidad, constituye la técnica estándar en programación competitiva para factorizar números de hasta 64 bits.

\graybold{Complejidad:}

Miller-Rabin determinístico para enteros de 64 bits: \bigO{\log^3 n}.

Pollard's Rho: tiempo esperado aproximado \bigO{n^{1/4}}.

Espacio \bigO{1}, sin contar la recursión utilizada para descomponer los factores.

\graybold{Fundamento teórico:}

Pollard's Rho genera una sucesión pseudoaleatoria
\[
x_{i+1}=f(x_i)\pmod n,
\]

normalmente utilizando
\[
f(x)=x^2+c.
\]

Si \(n\) es compuesto, la secuencia eventualmente forma ciclos distintos módulo cada uno de sus factores primos. Utilizando el algoritmo de Floyd (tortuga y liebre), se detecta cuándo dos valores coinciden en uno de esos factores.

En cada iteración se calcula
\[
d=\gcd(|x-y|,n).
\]

Si
\[
1<d<n,
\]

entonces se ha encontrado un factor no trivial del número.

Antes de ejecutar Pollard's Rho se utiliza Miller-Rabin para determinar rápidamente si el número ya es primo.

\graybold{Ejercicio aplicado}

Dado un entero positivo \(n\), imprimir su factorización prima ordenada de menor a mayor.

Los números pueden alcanzar valores de hasta \(10^{18}\), por lo que una factorización mediante división por prueba resulta inviable.

\graybold{I/O:}

Se reciben múltiples casos de prueba.

Cada caso contiene un único entero de 64 bits.

La implementación utiliza Miller-Rabin determinístico para enteros de 64 bits junto con Pollard's Rho para obtener todos los factores primos.

\graybold{Código solución}

\begin{lstlisting}[language=Python]
import random
from math import gcd

def is_prime(n):
    if n < 2:
        return False
    small = [2,3,5,7,11,13,17,19,23,29]

    for p in small:
        if n % p == 0:
            return n == p

    d = n - 1
    s = 0

    while d % 2 == 0:
        s += 1
        d //= 2

    for a in [2,325,9375,28178,450775,9780504,1795265022]:
        if a % n == 0:
            continue

        x = pow(a, d, n)

        if x == 1 or x == n - 1:
            continue

        for _ in range(s - 1):
            x = pow(x, 2, n)

            if x == n - 1:
                break
        else:
            return False

    return True


def pollard(n):
    if n % 2 == 0:
        return 2

    while True:
        c = random.randrange(1, n)
        f = lambda x: (pow(x, 2, n) + c) % n

        x = random.randrange(0, n)
        y = x
        d = 1

        while d == 1:
            x = f(x)
            y = f(f(y))
            d = gcd(abs(x - y), n)

        if d != n:
            return d


def factor(n, ans):
    if n == 1:
        return

    if is_prime(n):
        ans.append(n)
        return

    d = pollard(n)
    factor(d, ans)
    factor(n // d, ans)


n = int(input())
factors = []
factor(n, factors)
factors.sort()
print(*factors)
\end{lstlisting}

\graybold{Observación:}

Pollard's Rho es un algoritmo probabilístico. Distintas ejecuciones pueden seguir recorridos diferentes debido a la elección aleatoria de la función \(f(x)\), pero todas producen la misma factorización prima. En programación competitiva es la técnica estándar para factorizar enteros de 64 bits y suele combinarse con Miller-Rabin para verificar primalidad antes de intentar encontrar un factor.